\section*{第 \ref{ch:b1:water-small-molecules} 章——水与小分子生物分子}

\begin{solution}{exo:b1:water-small-molecules:1}
弯曲的\omterm{def:b1:water-small-molecules:water}{极性分子}，氧原子带部分负电荷，
氢原子带部分正电荷：以氧原子包围阳离子，
以氢原子包围阴离子，使其稳定（水合），
并与极性基团形成\omterm{def:b1:water-small-molecules:water}{氢键}。油没有可与之成键的电荷
或极性基团；油周围的水分子必须排列成笼状，这是不利的，
因此油被排斥。
\end{solution}

\begin{solution}{exo:b1:water-small-molecules:2}
\omterm{def:b1:water-small-molecules:water}{氢键}（\qtyrange{4}{20}{kJ/mol}）：DNA 中的碱基配对、蛋白质
二级结构。\omterm{def:b1:water-small-molecules:bonds}{离子键}（水中约 \qty{12}{kJ/mol}）：盐
桥（蛋白质中）、底物结合。\omterm{prop:b1:water-small-molecules:properties}{疏水效应}（不是化学键，
但总体能量相当）：膜的组装、蛋白质折叠。范德瓦耳斯力
（每个作用 \qtyrange{1}{4}{kJ/mol}）：任意两个互补表面的紧密堆积、
酶与底物的契合。
\end{solution}

\begin{solution}{exo:b1:water-small-molecules:3}
$\mathrm{pH} = -\log(4\times 10^{-8}) = 7.4$。在 \omterm{def:b1:water-small-molecules:ph}{pH} 1.5 时，
$[\mathrm{H^+}] = 10^{-1.5} = \qty{0.032}{mol/L}$。
\end{solution}

\begin{solution}{exo:b1:water-small-molecules:4}
约从 5.8 到 7.8（p$K_a \pm 1$）。在 7.4 时：$\log(\text{碱}/\text{酸})
= 0.6$，比值为 $4$。
\end{solution}

\begin{solution}{exo:b1:water-small-molecules:5}
在 \omterm{def:b1:water-small-molecules:ph}{pH} 4.76 时：\qty{50}{\%}。在 7.2 时：$\log(\mathrm{A^-}/\mathrm{HA}) =
2.44$，比值为 $275$：\qty{99.6}{\%} 发生解离。在 \omterm{def:b1:water-small-molecules:ph}{pH} 2 时，比值为 $10^{-2.76} =
1/575$：\qty{0.2}{\%} 发生解离。中性的酸形式可穿过\omterm{def:b1:membranes-transport:membrane}{脂
双层}；乙酸在胃中被吸收，此时它未解离，
进入\omterm{def:b1:eukaryotic-cell:organelle}{细胞溶胶}后以乙酸根形式被截留。
\end{solution}

\begin{solution}{exo:b1:water-small-molecules:6}
\omterm{def:b1:water-small-molecules:ph}{pH} $= 6.8 + \log 1 = 6.8$。加入盐酸后，碱为 $40$，酸为 $60$：\omterm{def:b1:water-small-molecules:ph}{pH} $= 6.8 +
\log(40/60) = 6.62$。在纯水中，$[\mathrm{H^+}] = \qty{0.01}{mol/L}$：
\omterm{def:b1:water-small-molecules:ph}{pH} $2$。
\end{solution}

\begin{solution}{exo:b1:water-small-molecules:7}
$500/2.4 = \qty{208}{g}$。要使一个分子汽化，它的四个
\omterm{def:b1:water-small-molecules:water}{氢键}都必须断裂；\qty{2.4}{kJ/g} 即 \qty{43}{kJ/mol}，大致相当于每个分子
两个键的能量（每个键由两个分子共用）。
\end{solution}

\begin{solution}{exo:b1:water-small-molecules:8}
$\Delta G^{\circ\prime} = -RT\ln 19 = -2.48\times 2.94 =
\qty{-7.3}{kJ/mol}$。在\omterm{def:b1:cell-unit-of-life:cell}{细胞}中：$\Delta G = -7.3 + 2.48\ln(2/0.02) =
-7.3 + 11.4 = \qty{+4.1}{kJ/mol}$：反应向逆方向进行，朝向
葡萄糖-1-磷酸（糖原合成中即是如此）。
\end{solution}

\begin{solution}{exo:b1:water-small-molecules:9}
水在 \qty{4}{\celsius} 时密度最大；较冷的水和冰密度更小，
留在表层，因此冰在上方形成，起隔热作用，
使下方液体整个冬季都保持在 \qty{4}{\celsius} 附近。若冰
下沉，湖泊将自底部向上冻结直至冻实，夏季只在
表层融化；淡水生物要熬过冬季
则是不可能的。
\end{solution}

\begin{solution}{exo:b1:water-small-molecules:10}
$RT = \qty{2.58}{kJ/mol}$。静息时：$Q = 0.9\times 10^{-3}\times 8\times
10^{-3}/8\times 10^{-3} = 9\times 10^{-4}$，$\ln Q = -7.0$：$\Delta G =
-30.5 - 18.1 = \qty{-48.6}{kJ/mol}$。疲劳时：$Q = 3\times 10^{-3}\times
30\times 10^{-3}/5\times 10^{-3} = 0.018$，$\ln Q = -4.0$：$\Delta G =
-30.5 - 10.4 = \qty{-40.9}{kJ/mol}$。疲劳使每个 \omterm{def:b1:water-small-molecules:atp}{ATP} 释放的能量减少了
\qty{16}{\%}，足以减慢那些耗能最多的
泵和马达。
\end{solution}

\begin{solution}{exo:b1:water-small-molecules:11}
$[\mathrm{H^+}] = 10^{-7.2} = \qty{6.3e-8}{mol/L}$；体积
\qty{1e-12}{L}：$\qty{6.3e-20}{mol}\times\num{6e23} = 38$ 个质子。相对于
$\num{3e9}$ 个缓冲基团：“\omterm{def:b1:water-small-molecules:ph}{pH}”并不是质子的计数，而是
\omterm{def:b1:water-small-molecules:buffer}{缓冲液}的质子化状态；缓冲基团与
水交换质子的速度比少量游离质子快十亿倍，后者可以忽略。
\end{solution}

\begin{solution}{exo:b1:water-small-molecules:12}
在平衡时，\omterm{def:b1:water-small-molecules:atp}{ATP}、ADP 与磷酸的比值约为
$10^{-5}$；\omterm{def:b1:cell-unit-of-life:cell}{细胞}中 \omterm{def:b1:water-small-molecules:atp}{ATP} 的浓度比 ADP 高出一千倍，因此其
水解释放 \qty{50}{kJ/mol}，而不是零。每一个吸能过程
——合成、运输、运动——都与
这一位移\omterm{def:b1:water-small-molecules:atp}{偶联}；分解代谢的唯一作用是维持它。若\omterm{def:b1:cell-unit-of-life:cell}{细胞}的 \omterm{def:b1:water-small-molecules:atp}{ATP}
达到平衡，则没有梯度可供消耗：泵停止工作，离子外漏，
\omterm{def:b1:cell-unit-of-life:cell}{细胞}肿胀，数分钟内死亡。活着，就是维持
这种非平衡状态，而非某种物质。
\end{solution}

\begin{solution}{pb:b1:water-small-molecules:1}
\textbf{1.} $6.1 + \log(24/1.2) = 6.1 + 1.301 = 7.40$。
\textbf{2.} $10^{-7.4} = \qty{4.0e-8}{mol/L} = \qty{40}{nmol/L}$。
\textbf{3.} $40\times 5 = \qty{200}{nmol}$，比碳酸氢根少
六十万倍。
\textbf{4.} $\mathrm{CO_2} = 1.2\times 46/40 = \qty{1.38}{mmol/L}$；
\omterm{def:b1:water-small-molecules:ph}{pH} $= 6.1 + \log(25/1.38) = 6.1 + 1.258 = 7.36$。
\textbf{5.} 在烧瓶中，$20:1$ 的比值使可供吸收碱的酸组分很少，
难以吸收碱，且在距 p$K_a$ 一个单位处缓冲容量很低；而在
人体内，酸组分是一种气体，其浓度由肺维持，
并可无限再生，因此比值可通过
通气重新设定，与 p$K_a$ 无关。
\textbf{6.} $24/1.2 = 20$；碱占该缓冲对的 $20/21 = \qty{95}{\%}$。
\textbf{7.} $60/5 = \qty{12}{mmol/L}$。
\textbf{8.} 碳酸氢根 $24 - 12 = \qty{12}{mmol/L}$；$\mathrm{CO_2}$
为 $1.2 + 12 = \qty{13.2}{mmol/L}$。
\textbf{9.} \omterm{def:b1:water-small-molecules:ph}{pH} $= 6.1 + \log(12/13.2) = 6.1 - 0.04 = 6.06$。
\textbf{10.} $[\mathrm{H^+}] = \qty{0.012}{mol/L}$：\omterm{def:b1:water-small-molecules:ph}{pH} $1.92$。
\textbf{11.} 变化了 $1.34$ 个单位：它使血液避免降至 \omterm{def:b1:water-small-molecules:ph}{pH} 1.9，
但 6.06 远低于任何酶所能耐受的范围。作为封闭
系统，效果很差。
\textbf{12.} $[\mathrm{H^+}]$ 从 \qty{40}{nmol/L} 升至 $10^{-6.06} = \qty{870}{nmol/L}$：加入 \qty{12}{mmol/L} 的酸使其上升 \qty{0.83}{\micro mol/L}，即每 \num{14000} 个质子中约有一个保持游离，其余都结合在碳酸氢根上。
\textbf{13.} \omterm{def:b1:water-small-molecules:ph}{pH} $= 6.1 + \log(12/1.2) = 6.1 + 1.0 = 7.10$。
\textbf{14.} 生成的 $\mathrm{CO_2}$：$\qty{12}{mmol/L}\times
\qty{5}{L} = \qty{60}{mmol}$。
\textbf{15.} 一分钟内产生 \qty{60}{mmol}，加上原有的 \qty{10}{mmol}：
增至原来的七倍。
\textbf{16.} $\mathrm{CO_2} = \qty{0.9}{mmol/L}$；\omterm{def:b1:water-small-molecules:ph}{pH} $= 6.1 +
\log(12/0.9) = 6.1 + 1.125 = 7.22$。
\textbf{17.} 在 7.4 附近的开放碳酸氢盐体系中：碳酸氢根降至十分之一时，\omterm{def:b1:water-small-molecules:ph}{pH} 改变一个单位，
即从 \qty{24}{} 降至
\qty{2.4}{mmol/L}，即每个单位消耗 \qty{21.6}{mmol/L} 的酸；但在
最初的几个单位内，局部斜率为 $2.3\times[\mathrm{HCO_3^-}] =
\qty{55}{mmol/L}$ 每单位。与蛋白质的 \qty{25}{mmol/L}
每单位相比，碳酸氢根结合约 $55/80 = \qty{69}{\%}$ 的质子，
蛋白质结合 \qty{31}{\%}。
\textbf{18.} 碳酸氢根结合 $0.69\times 12 = \qty{8.3}{mmol/L}$：
$24 - 8.3 = \qty{15.7}{mmol/L}$；\omterm{def:b1:water-small-molecules:ph}{pH} $= 6.1 + \log(15.7/1.2) = 6.1 +
1.117 = 7.22$（蛋白质份额与 \qty{25}{mmol/L}
每单位 $\times 0.18$ 单位 $= \qty{4.5}{mmol/L}$ 相符，接近其
\qty{3.7}{mmol/L}）。
\textbf{19.} 去除酸后，被消耗的碳酸氢根得以再生
（从 \qty{12}{mmol/L} 恢复至 \qty{24}{}），在 $\mathrm{CO_2}$ 保持不变时，
\omterm{def:b1:water-small-molecules:ph}{pH} 回到 7.40。
\textbf{20.} $60/5 = \qty{12}{mmol/L}$ 的碳酸氢根每天流失；
三天后 $24 - 36 < 0$：碳酸氢根耗尽，\omterm{def:b1:water-small-molecules:ph}{pH}
大幅降至 7 以下——实际中患者必须补充
碳酸氢盐或进行透析；仅过一天时，$6.1 + \log(12/1.2) =
7.10$。
\textbf{21.} $\mathrm{CO_2} = 1.2\times 60/40 = \qty{1.8}{mmol/L}$；
\omterm{def:b1:water-small-molecules:ph}{pH} $= 6.1 + \log(24/1.8) = 6.1 + 1.125 = 7.22$。若碳酸氢根为
\qty{33}{}：$6.1 + \log(33/1.8) = 6.1 + 1.263 = 7.36$。
\textbf{22.} 肺改变分母（溶解的
$\mathrm{CO_2}$），因其为肺所呼出的气体，可在数次呼吸内改变；
肾脏改变分子（碳酸氢根），通过排出酸并
使碳酸氢根再生来实现，这一过程需要数小时至数天。
\textbf{23.} 在 \omterm{def:b1:water-small-molecules:ph}{pH} 5 时，$[\mathrm{H^+}] = \qty{10}{\micro mol/L}$：\qty{1.5}{L} 中游离的为 \qty{15}{\micro mol}，仅为 \qty{60}{mmol} 的四千分之一；酸以与缓冲物质结合的形式排出——即铵和磷酸二氢根。
\textbf{24.} 失去酸会使 \omterm{def:b1:water-small-molecules:ph}{pH} 升高：碳酸氢根增加
$100/5 = \qty{20}{mmol/L}$，升至 \qty{44}{}；\omterm{def:b1:water-small-molecules:ph}{pH} $= 6.1 + \log(44/1.2) =
6.1 + 1.564 = 7.66$（碱中毒），在肺和肾
代偿之前。
\textbf{25.} 封闭系统：由 7.40 降至 6.06，变化 $-1.34$。开放系统，
仅含碳酸氢盐：降至 7.10，变化 $-0.30$。开放系统，含蛋白质：
降至 7.22，变化 $-0.18$。
\end{solution}
